% Part: incompleteness
% Chapter: representability-in-q
% Section: introduction

\documentclass[../../../include/open-logic-section]{subfiles}

\begin{document}

\olfileid{inc}{req}{int}
\olsection{Pambuka}

Teorema-teorema ora-jangkep lumaku kanggo teori kang bisa
nyatakake lan mbuktekake fakta dhasar ngenani fungsi komputabel.
Kita bakal njlentrehake teori kang minimal banget kaya mangkono,
diarani ``$\Th{Q}$'' (utawa kadhangkala ``$Q$ Robinson,''
miturut Raphael Robinson). Kita bakal nerangake teges fungsi
kang \emph{representabel} ing $\Th{Q}$, banjur mbuktekake iki:
\begin{quote}
  Sawijining fungsi representabel ing $\Th{Q}$ yen lan mung yen komputabel.
\end{quote}
Iki menehi model komputabilitas liyane. Kajaba iku, kita bakal
nggunakake kanggo nuduhake yen himpunan
$\Setabs{!A}{\Th{Q} \Proves !A}$ ora bisa diputusake, kanthi
ngreduksi masalah mandheg marang himpunan iki. Ing pungkasan,
kita bakal mbuktekake pratelan kang luwih kuwat tinimbang iki.

Basa $\Th{Q}$ iku basa aritmetika; $\Th{Q}$ dumadi saka
aksioma-aksioma iki (kang digunakake bebarengan karo aksioma
lan aturan liyane saka logika orde kapisan kanthi !!{identity}):
\begin{align*}
& \lforall[x][\lforall[y][(\eq[x'][y'] \lif \eq[x][y])]] \tag{$!Q_1$}\\
& \lforall[x][\eq/[\Obj 0][x']] \tag{$!Q_2$}\\
& \lforall[x][(\eq[x][\Obj 0] \lor \lexists[y][\eq[x][y']])] \tag{$!Q_3$}\\
& \lforall[x][\eq[(x + \Obj 0)][x]] \tag{$!Q_4$}\\
& \lforall[x][\lforall[y][\eq[(x + y')][(x + y)']]] \tag{$!Q_5$}\\
& \lforall[x][\eq[(x \times \Obj 0)][\Obj 0]] \tag{$!Q_6$}\\
& \lforall[x][\lforall[y][\eq[(x \times y')][((x \times y) + x)]]] \tag{$!Q_7$}\\
& \lforall[x][\lforall[y][(x < y \liff \lexists[z][\eq[(z' + x)][y]])]] \tag{$!Q_8$}
\end{align*}
Kanggo saben bilangan asli $n$, tetepna numeral $\num{n}$ minangka
suku $\Obj{0}^{\prime\prime\ldots\prime}$ kang nduweni $n$
tandha prima kabeh. Mula $\num{0}$ iku !!{constant}~$\Obj{0}$
wae, $\num{1}$ iku $\Obj{0}'$, $\num{2}$ iku $\Obj{0}''$,
lan sapiturute.

Minangka teori aritmetika, $\Th{Q}$ iku \emph{ringkih banget};
umpamane, fakta kang prasaja kaya
$\lforall[x][\eq/[x][x']]$ utawa $\lforall[x][\lforall[y][(x + y) = (y +
    x)]]$ wae ora bisa dibuktekake. Nanging kita bakal weruh yen
$\Th{Q}$ narik kawigaten malah \emph{amarga} ringkihe kuwi.
Sejatine, kekuwatane mung pas cukup supaya teorema ora-jangkep
lumaku. Alesan liyane kenapa $\Th{Q}$ narik kawigaten yaiku
himpunan aksiomane \emph{winates}.

Teori kang luwih kuwat tinimbang $\Th{Q}$ (diarani
\emph{aritmetika Peano} $\Th{PA}$) diduweni kanthi nambahake
skema induksi marang~$\Th{Q}$:
\[
(!A(\Obj 0) \land \lforall[x][(!A(x) \lif !A(x'))]) \lif \lforall[x][!A(x)]
\]
ing ngendi $!A(x)$ iku formula sembarang. Yen $!A(x)$ ngemot
!!{variable}s bebas liyane saliyane $x$, kita nambahake
kuantor universal ing ngarep kanggo ngiket kabeh variabel mau
(supaya instans skema induksi kang cocog iku !!a{sentence}).
Upamane, yen $!A(x, y)$ uga ngemot !!{variable}~$y$ kang
bebas, instans kang cocog yaiku
\[
\lforall[y][((!A(\Obj 0) \land \lforall[x][(!A(x) \lif !A(x'))]) \lif
  \lforall[x][!A(x)])]
\]
Kanthi nggunakake instans skema induksi, luwih akeh pratelan
bisa dibuktekake saka aksioma~$\Th{PA}$ tinimbang saka
aksioma $\Th{Q}$. Malah, golek pratelan ``alamiah'' ngenani
bilangan asli kang ora bisa dibuktekake ing aritmetika Peano
mbutuhake gaweyan kang akeh!{}

\begin{defn}
\ollabel{defn:representable-fn}
  Fungsi $f(x_0,\ldots,x_k)$ saka bilangan asli menyang
  bilangan asli diarani {\em representabel ing $\Th{Q}$} yen
  ana formula $!A_f(x_0,\dots,x_k,y)$ saengga saben
  $f(n_0,\dots,n_k) = m$, $\Th{Q}$ mbuktekake
\begin{enumerate}
\item\ollabel{defn:rep:a} $!A_f(\num{n_0}, \dots, \num{n_k}, \num{m})$
\item\ollabel{defn:rep:b} $\lforall[y][(!A_f(\num{n_0}, \dots,
\num{n_k}, y) \lif \num{m} = y)]$.
\end{enumerate}
\end{defn}

Ana cara liyane kanggo nyatakake definisi iki; upamane,
kanthi padha tegese kita bisa mbutuhake supaya $\Th{Q}$
mbuktekake $\lforall[y][(!A_f(\num{n_0}, \dots,
\num{n_k}, y) \liff \eq[y][\num{m}])]$.

\begin{thm}
\ollabel{thm:representable-iff-comp}
Sawijining fungsi representabel ing $\Th{Q}$ yen lan mung yen komputabel.
\end{thm}

Ana rong arah kanggo mbuktekake teorema iki. Arah saka kiwa
menyang tengen cukup gampang sawisé sintaks wis diaritmetisasi.
Arah sijine mbutuhake gaweyan luwih akeh. Gagasan dhasare:
kita nggunakake ``rekursif umum'' kanggo menehi teges kang
pas marang ``komputabel'', banjur nuduhake yen saben fungsi
rekursif umum representabel ing~$\Th{Q}$. Elinga yen fungsi
iku rekursif umum yen bisa ditegesi saka $\Zero$, fungsi
penerus~$\Succ$, lan fungsi proyeksi~$\Proj{n}{i}$, kanthi
komposisi, rekursi primitif, lan panggolekan tanpa wates reguler.
Mula salah siji cara nuduhake yen saben fungsi rekursif umum
representabel ing~$\Th{Q}$ yaiku nuduhake yen fungsi dhasar
representabel, lan yen sawetara fungsi representabel, mula
fungsi kang ditegesi saka fungsi mau nganggo komposisi,
rekursi primitif, lan panggolekan tanpa wates reguler uga
representabel. Kanthi tembung liya, kita bisa nuduhake yen
fungsi dhasar representabel lan himpunan fungsi representabel
``katutup tumrap'' komposisi, rekursi primitif, lan panggolekan
tanpa wates reguler. Iki njamin saben fungsi rekursif umum
representabel.

Nyatane, langkah kanggo nuduhake yen fungsi representabel
katutup tumrap rekursi primitif iku angel. Supaya ora perlu
langkah iki, kita luwih dhisik nuduhake yen rekursi primitif
sejatine bisa ditinggalake. Tegese, saben fungsi rekursif umum
bisa ditegesi saka fungsi dhasar mung nganggo komposisi lan
panggolekan tanpa wates reguler. Kanggo iki, kita nuduhake
yen rekursi primitif bisa ditindakake liwat sawijining
panggolekan tanpa wates reguler kang tartamtu. Nanging supaya
iki bisa lumaku, kita kudu nambah sawetara fungsi dhasar:
panjumlahan, perkalian, lan fungsi karakteristik saka relasi
identitas~$\Char{=}$. Banjur kita bisa mbuktekake teorema
kanthi nuduhake yen kabeh fungsi dhasar \emph{iki}
representabel ing~$\Th{Q}$, lan fungsi representabel katutup
tumrap komposisi lan panggolekan tanpa wates reguler.

\end{document}
